\documentclass[12pt,a4paper,oneside,pdftex]{report}

\usepackage[latin1]{inputenc}
\usepackage[OT1]{fontenc}
\usepackage[finnish,swedish,english]{babel}
% \usepackage{palatino}
% \usepackage{tgpagella}


% Optional packages
% ------------------------------------------------------------------

% Natbib allows you to select the format of the bibliography references.
% The first example uses numbered citations: 
\usepackage[square,sort&compress,numbers]{natbib}
% Also, you should use \citet (cite in text) when you wish to refer to the author directly (\citet{blaablaa} said blaa blaa), and \citep when you wish to refer similarly than with numbered citations
% (It has been said that blaa blaa~\citep{blaablaa}).

% The eurosym package provides a euro symbol. Use with \euro{}
%\usepackage{eurosym} 

% The listing package provides automatic code formatting utilities so that you can copy-paste code examples and have them rendered nicely.
% See the package documentation for details.
\usepackage{listings}

% Supertabular provides a tabular environment that can span multiple pages. 
%\usepackage{supertabular}
% Longtable provides a tabular environment that can span multiple pages.
% This is used in the example acronyms file. 
\usepackage{longtable}

% The fancyhdr package allows you to set your the page headers manually, and allows you to add separator lines and so on. 
% Check the package documentation. 
% \usepackage{fancyhdr}

% Subfigure package allows you to use subfigures (i.e. many subfigures within one figure environment).
% These can have different labels and they are numbered automatically.
% Check the package documentation. 
\usepackage{subfigure}

% The titlesec package can be used to alter the look of the titles of sections, chapters, and so on.
% This example uses the ``medium'' package option which sets the titles to a medium size, making them a bit smaller than what is the default.
% You can fine-tune the title fonts and sizes by using the package options.
% See the package documentation.
\usepackage[medium]{titlesec}

% The TikZ package allows you to create professional technical figures.
% The learning curve is quite steep, but it is definitely worth it if  you wish to have really good-looking technical figures. 
\usepackage{tikz}
% You also need to specify which TikZ libraries you use
\usetikzlibrary{positioning}
\usetikzlibrary{calc}
\usetikzlibrary{arrows}
\usetikzlibrary{decorations.pathmorphing,decorations.markings}
\usetikzlibrary{shapes}
\usetikzlibrary{patterns}


% Aalto-thesis
% ------------------------------------------------------------------

% The aalto-thesis package provides typesetting instructions for the standard master's thesis parts (abstracts, front page, and so on).
% Load this package second-to-last, just before the hyperref package.
% Options that you can use: 
%   mydraft - renders the thesis in draft mode. 
%             Do not use for the final version. 
%   doublenumbering - [optional] number the first pages of the thesis
%                     with roman numerals (i, ii, iii, ...); and start
%                     arabic numbering (1, 2, 3, ...) only on the 
%                     first page of the first chapter
%   twoinstructors  - changes the title of instructors to plural form
%   twosupervisors  - changes the title of supervisors to plural form
\usepackage[mydraft,twosupervisors]{aalto-thesis}
%\usepackage[mydraft,doublenumbering]{aalto-thesis}
%\usepackage{aalto-thesis}


% Hyperref
% ------------------------------------------------------------------
% Hyperref creates links from URLs, for references, and creates a TOC in the PDF file.
% This package must be the last one you include, because it has compatibility issues with many other packages and it fixes those issues when it is loaded.   
\RequirePackage[pdftex]{hyperref}
% Setup hyperref so that links are clickable but do not look different
\hypersetup{colorlinks=false,raiselinks=false,breaklinks=true}
\hypersetup{pdfborder={0 0 0}}
\hypersetup{bookmarksnumbered=true}
% The following line suggests the PDF reader that it should show the first level of bookmarks opened in the hierarchical bookmark view. 
\hypersetup{bookmarksopen=true,bookmarksopenlevel=1}
% Hyperref can also set up the PDF metadata fields. 
% These are set a bit later on, after the thesis setup.   


% Thesis setup
% ==================================================================
% Change these to fit your own thesis.
% \COMMAND always refers to the English version;
% \FCOMMAND refers to the Finnish version; and
% \SCOMMAND refers to the Swedish version.
% You may comment/remove those language variants that you do not use (but then you must not include the abstracts for that language)
% ------------------------------------------------------------------
% If you do not find the command for a text that is shown in the cover page or in the abstract texts, check the aalto-thesis.sty file and locate the text from there. 
% All the texts are configured in language-specific blocks (lots of commands that look like this: \renewcommand{\ATCITY}{Espoo}.
% You can just fix the texts there.
% Just remember to check all the language variants you use (they are all there in the same place). 
% ------------------------------------------------------------------
\newcommand{\TITLE}{English Name PLACE-HOLDER}
\newcommand{\FTITLE}{Finnish Name PLACE-HOLDER}
\newcommand{\SUBTITLE}{English Subtitle PLACE-HOLDER}
\newcommand{\FSUBTITLE}{Finnish Subtitle PLACE-HOLDER}
\newcommand{\DATE}{English Date PLACE-HOLDER}
\newcommand{\FDATE}{Finnish Date PLACE-HOLDER}


% Supervisors and instructors
% ------------------------------------------------------------------
% Example of one supervisor:
\newcommand{\SUPERVISOR}{English Supervisor PLACE-HOLDER}
\newcommand{\FSUPERVISOR}{Finnish Supervisor PLACE-HOLDER}

% If you have only one instructor, just write one name here
\newcommand{\INSTRUCTOR}{English Instructor PLACE-HOLDER}
\newcommand{\FINSTRUCTOR}{Finnish Instructor PLACE-HOLDER}


% Other stuff
% ------------------------------------------------------------------
\newcommand{\PROFESSORSHIP}{English Professorship PLACE-HOLDER}
\newcommand{\FPROFESSORSHIP}{Finnish Professorship PLACE-HOLDER}
% Professorship code is the same in all languages
\newcommand{\PROFCODE}{Code PLACE-HOLDER}
\newcommand{\KEYWORDS}{English Keywords PLACE-HOLDER}
\newcommand{\FKEYWORDS}{Finnish Keywords PLACE-HOLDER}
\newcommand{\LANGUAGE}{English}
\newcommand{\FLANGUAGE}{Englanti}

% Author is the same for all languages
\newcommand{\AUTHOR}{Mari Nikkarinen}

% Currently the English versions are used for the PDF file metadata
% Set the PDF title
\hypersetup{pdftitle={\TITLE\ \SUBTITLE}}
% Set the PDF author
\hypersetup{pdfauthor={\AUTHOR}}
% Set the PDF keywords
\hypersetup{pdfkeywords={\KEYWORDS}}
% Set the PDF subject
\hypersetup{pdfsubject={Master's Thesis}}


% Layout settings
% ------------------------------------------------------------------

% If you write your thesis Finnish, uncomment these lines; if you write in English, leave these lines commented! 
% \setlength{\parindent}{0pt}
% \setlength{\parskip}{1ex}

% Use this to control how much space there is between each line of text.
% 1 is normal (no extra space), 1.3 is about one-half more space, and 1.6 is about double line spacing.  
% \linespread{1} % This is the default
% \linespread{1.3}

% Bibliography style
% acm style gives you a basic reference style.
% It works only with numbered references.
%\bibliographystyle{acm}
% Plainnat is a plain style that works with both numbered and name citations.
\bibliographystyle{plainnat}


% Extra hyphenation settings
% ------------------------------------------------------------------
% You can list here all the files that are not hyphenated correctly.
% You can provide many \hyphenation commands and/or separate each word with a space inside a single command.
% Put hyphens in the places where a word can be hyphenated.
% Note that (by default) LaTeX will not hyphenate words that already have a hyphen in them (for example, if you write ``structure-modification operation'', the word structure-modification will never be hyphenated).
% You need a special package to hyphenate those words.
\hyphenation{di-gi-taa-li-sta yksi-suun-tai-sta}



% The preamble ends here, and the document begins. 
% Place all formatting commands and such before this line.
% ------------------------------------------------------------------
\begin{document}
% This command adds a PDF bookmark to the cover page.
% You may leave it out if you don't like it...
\pdfbookmark[0]{Cover page}{bookmark.0.cover}
% This command is defined in aalto-thesis.sty.
% It controls the page numbering based on whether the doublenumbering option is specified
\startcoverpage


% Cover page
% ------------------------------------------------------------------
% Options: finnish, english, and swedish
% These control in which language the cover-page information is shown
\coverpage{english}

% Abstra
cts
% ------------------------------------------------------------------
% Include an abstract in the language that the thesis is written in, and if your native language is Finnish or Swedish, one in that language.


% Abstract in English
% ------------------------------------------------------------------
\thesisabstract{english}{

\fixme{The abstract is written last.}}

% Abstract in Finnish
% ------------------------------------------------------------------
\thesisabstract{finnish}{
\fixme{Abstrakti kirjoitetaan viimeisen�.}}

% Acknowledgements
% ------------------------------------------------------------------
% Select the language you use in your acknowledgements
\selectlanguage{english}

% Uncomment this line if you wish acknoledgements to appear in the table of contents
%\addcontentsline{toc}{chapter}{Acknowledgements}

% The star means that the chapter isn't numbered and does not show up in the TOC
\chapter*{Acknowledgements}

\fixme{Acknowledge some people.}

\vskip 10mm

\noindent Espoo, \DATE
\vskip 5mm
\noindent\AUTHOR

% Acronyms
% ------------------------------------------------------------------
% Use \cleardoublepage so that IF two-sided printing is used (which is not often for masters theses), then the pages will still start correctly on the right-hand side.
\cleardoublepage
% Example acronyms are placed in a separate file, acronyms.tex
\addcontentsline{toc}{chapter}{Abbreviations and Acronyms}
\chapter*{Abbreviations and Acronyms}

% The longtable environment should break the table properly to multiple pages, 
% if needed

\noindent
\begin{longtable}{@{}p{0.25\textwidth}p{0.7\textwidth}@{}}
2k/4k/8k mode & COFDM operation modes \\
3GPP & 3rd Generation Partnership Project \\ 
ESP & Encapsulating Security Payload; An IPsec security protocol \\ 
FLUTE  & The File Delivery over Unidirectional Transport protocol \\ 
e.g.& for example (do not list here this kind of common acronymbs or abbreviations, but only those that are essential for understanding the content of your thesis. \\ 
note & Note also, that this list is not compulsory, and should be omitted if you have only few abbreviations

\end{longtable}


% Table of contents
% ------------------------------------------------------------------
\cleardoublepage
% This command adds a PDF bookmark that links to the contents.
% You can use \addcontentsline{} as well, but that also adds contents entry to the table of contents, which is kind of redundant.
% The text ``Contents'' is shown in the PDF bookmark. 
\pdfbookmark[0]{Contents}{bookmark.0.contents}
\tableofcontents

% List of tables
% ------------------------------------------------------------------
% You only need a list of tables for your thesis if you have very many tables.
% If you do, uncomment the following two lines.
% \cleardoublepage
% \listoftables

% Table of figures
% ------------------------------------------------------------------
% You only need a list of figures for your thesis if you have very many figures.
% If you do, uncomment the following two lines.
% \cleardoublepage
% \listoffigures

% The following label is used for counting the prelude pages
\label{pages-prelude}
\cleardoublepage

%%%%%%%%%%%%%%%%% The main content starts here %%%%%%%%%%%%%%%%%%%%%
% ------------------------------------------------------------------
% This command is defined in aalto-thesis.sty.
% It controls the page numbering based on whether the doublenumbering option is specified
\startfirstchapter

% Add headings to pages (the chapter title is shown)
\pagestyle{headings}

% The contents of the thesis are separated to their own files.
% Edit the content in these files, rename them as necessary.
% ------------------------------------------------------------------
\chapter{Introduction}
\label{chapter:intro}

This is my master's thesis, and I am very proud of it.  Of course,
when I write my \emph{real} master's thesis, I will not use the
singular pronoun \emph{I}, but rather try to avoid referring to myself
and speak of the research \emph{we} have conducted---I rarely work
alone, after all.  Yet, both \emph{I} and \emph{we} are correct, and
it depends on the instructor and the supervisor (of course from you,
too), which one they would prefer. Anyway, the tense should be active,
and passive sentenses should be avoided (especially, writing sentences
where the subject is presented with by preposition), so often you
cannot avoid choosing between the pronouns. Life is strange, but there
you have it.

By the way, the preferred order of writing your master's thesis is
about the same as the outline of the thesis: you first discover your
problem and write about that, then you find out what methods you
should use and write about that.  Then you do your implementation, and
document that, and so on.  However, the abstract and introduction are
often easiest to write last.  This is because these really cover the
entire thesis, and there is no way you could know what to put in your
abstract before you have actually done your implementation and
evaluation. Rarely anyone write the thesis from the beginning to the
end just one time, but the writing is more like process, where every
piece of text is written at least twice. Be also prepared to delete
your own text. In the first phase, you can hide it into comments that
are started with \% but during the writing, the many comments should be
visible for your helpers, the instructor and supervisor.

The introduction in itself is rarely very long; two to five pages often
suffice.


\section{Problem statement}

Undergraduate students studying technical subjects do not consider typography
very interesting these days, and therefore the typographical quality of many
theses is unacceptably low. 
We plan to rectify this situation somewhat by providing a decent-quality
example thesis outline for students.
We expect that the typographical quality of the master's theses will
dramatically increase as the new thesis outline is taken into use.

\section{Helpful hints}

Read the information from the university master's thesis
pages~\cite{ThesisInstructions} before starting the thesis.  You
should also go through the thesis grading
instructions~\cite{ThesisGrading} together with your instructor and/or
supervisor in the beginning of your work.

\section{Structure of the Thesis}
\label{section:structure} 

You should use transition in your text, meaning that you should help
the reader follow the thesis outline. Here, you tell what will be in
each chapter of your thesis. 



\chapter{Background}
\label{chapter:background} 

The problem must have some background, otherwise it is not
interesting.  You can explain the background here. Probably you should
change the title to something that describes more the content of this
chapter. Background consists of information that help other masters of
the same degree program to understand the rest of the thesis.

Transitions mentioned in Section~\ref{section:structure} are used also
in the chapters and sections. For example, next in this chapter we
tell how to use English language, how to find and refer to sources,
and enlight different ways to include graphics in the thesis.

\section{Language and Structure}

Moreover, the transitions are also used in the paragraph and the
sentence level meaning that all the text is linked together. For example,
the word ``moreover'' here is one way, but of course you should use
variation in the text. Examples of transitional devices (words) and
their use can be found from writing guides, e.g. from the Online
Writing Lab
(OWL)\footnote{http://owl.english.purdue.edu/owl/resource/574/02/} of
Purdue University or Strunk's Elements of
Style\footnote{http://www.bartleby.com/141/}. Remember that footnotes
are additional information, and they are seldom used.  If you refer to a source, you do no
not use footnote. The right command for the references is \emph{cite}.

The language used in the thesis should be technical (or
scientifical). For example, the abreviations aren't used but all them
are written open (i.e. ``are not''). Since the content itself is often
hard to understand (and explain), the sentences should not be very
long, use complex language with several examples embedded in the same
sentence, and, also, seldom used words and weird euphemism or paraphrases
can make the sentence hard to follow and to read it with only one
time, and making everything even harder to understand all this without
any punctuation marks makes the instructor cry and finally after
trying to correct the language, you will get boomerang, and everyone's
time has just been wasted.

Please use proofreaders before sending even your unfinished version to
the instructor and/or supervisor. You will get better comments when
they do not need first proofread your text. Moreover, they can
consentrate to the content better if the language and spelling
mistakes are not distracting the reading. Several editors have their
own proofreading tools, e.g. ispell in emacs. You can also use
Microsoft Word to proofread your thesis: it can correct also some
grammatical errors and not just misspelled words.

Note also that if you have a section or a subsection, you have to have
at least two of them, or otherwise the section or subsection title is
unnecessary. Same with the paragraphs an: you should not have sections
with only one paragraphe, and single sentence paragraph. Furthermore,
always write some text after the title before the next level title.

\section{Finding and referring to sources}

Never ever copy anything into your theses from somebody else's text
(nor your own previously published text). Never. Not even for starting
point to be rewritten later. The risk is that you forgot the copied
text to your thesis and end up to be accused of plagiarism. Plagiarism
is a serious crime in studies and science and can ruin your career
even its beginning. To repeate: never cut and paste text into your
thesis!

\subsection{Finding sources}

All work is based on someone else's work. You should find the relevant
sources of your field and choose the best of them. Also, you should
refer to the original source where a fact has been mentioned first
time. Remember source evaluation (criticism) with all sources you
find.

Good starting points for finding references in computer science are: 
\begin{itemize}
% You can use this command to set the items in the list closer to each other
% (ITEM SEParation, the vertical space between the list items) 
\setlength{\itemsep}{0pt}
\item Nelli Portal (Aalto Library): \url{http://www.nelliportaali.fi}
\item ACM Digital library: \url{http://portal.acm.org/}
\item IEEExplore: \url{http://ieeexplore.ieee.org}
\item ScienceDirect: \url{http://www.sciencedirect.com/}
\item \ldots although Google Scholar (\url{http://scholar.google.com/}) will
find links to most of the articles from the abovementioned sources, if you
search from within the university network
\end{itemize}

Some of the publishers do not offer all the text of the articles
freely, but the library has agreed on the rights to use the whole
text. Thus, you should sometimes use computers in the domain of the
university in order to get the full text. Sometimes the Nelli Portal
can also help getting the whole article instead of just the abstract.
The library has also brief instrucions how to find
information~\cite{howfindinfo}.

Instead of normal Google, use Google Scholar
(\url{http://scholar.google.fi/}). It finds academic publications whereas
normal Google find too much commercial advertisements or otherwise
biased information. Wikipedia articles should be referred to in the master
thesis only very, very seldomly. You can use Wikipedia for understanding
some basics and finding more sources, but often you cannot be sure if
the article is correct and unbiased.

One important part of the sources that you have found is the reference
list. This way you can find the original sources that all the other
research of the field refer. Often you can also find more information
with the name of the researchers that are often referred in the
articles.

\subsection{Referring to sources}

The main point in referring to sources is to separate your own
thinking and text from that of others. Facts of the research area can
be given without reference, but otherwise you should refer to
sources. This means two things: marking the source in the text where
it has been used, and listing the sources usually in the end of the
thesis in a way that help the reader to find the original source.

There are several bibliography styles, meaning how to form the
bibliography in the end of the thesis. Aalto's library has good
instructions for many styles~\cite{bibinstructions}. You should ask
from your supervisor or instructors which style you should use. This
thesis template uses the number style that is often used in software
engineering. The other style also used in the CS field,
e.g. usability, is the Harvard style where instead of numbers, the
reference is marked into the text with author's name and publishing
year. Other areas use also many other styles for making the lists and
marking the references.

In addition to the list in the end of the thesis, you have to mark the
source in the text where the source is used. There are three places
for the reference: in a sentence before the period, in the end of a
sentence after the period, or in the end of a paragraph. All of them
have different meaning. The main point is that first you paraphrase
the source using your own words and then mark the source. Next, we
give short examples that are marked with \emph{emphasised text}.

\emph{Haapasalo~\cite{HaapasaloThesis} researched database algorithms
  that allows use of previous versions of the content stored in the
  database.} This kind of marking means that this paragraph (or until
the next reference is given) is based on the source mentioned in the
beginning.  Giving the source you should use only the family name of
the first author of the article, and not give any hints about what is
the type of the article that is referred.

\emph{B+-trees offers one way to index data that is stored in to a
  database. Multiversion B+-trees (MVBT) offer also a way to restore
  the data from previous versions of the database. Concurrent MVBT
  allows many simultaneous updates to the database that is was not
  possible with MVBT.~\cite{HaapasaloThesis}} When the marking is
after the period, the reference is retrospective: all the paragraph
(or after previous reference marking) is based on the source given in
its end. If the content is very broad, you can start with saying
\emph{According to Haapasalo}, then continue referring the source with
several separate sentences, and in the end put the marking of your
source \emph{ that shows that CMVBT are the
  best. ~\cite{HaapasaloThesis}}. 

If your paragraph has several sources, the above mentioned styles are
not proper. The reader of your thesis cannot know which of your
sources give which of the statements. In this case, it is better to
use more finegraded refering where the reference markings that are
embedded in the sentences. For example, \emph{the multiversion B+-tree
  (MVBT) index of Becker et al.~\cite{becker:1996:mvbt} allows database
  users to query old versions of the database, but the index is not
  transactional.
  It's successor, the transactional MBVT (TMVBT), allows a single transaction
  running in its own thread or process to update the database concurrently
  with other transactions that only read the
  database~\cite{haapasalo:2009:tmvbt}. 
  Further development, titled the concurrent MBVT (CMVBT),
  allows several transactions to perform updates to the database at the same
  time~\cite{HaapasaloThesis}}. 
  Here, the references are marked before
  the period in the sentences where they are used.

Finally, direct quotes are allowed. However, often you should avoid
them since they do not usually fit in to your text very well. Using
direct quotes has two tricks: quotation marks and the source.  \emph{
  ``Even though deletions in a multiversion index must not physically
  delete the history of the data items, queries and range scans can
  become more efficient, if the leaf pages of the index structure are
  merged to retain optimality.''~\cite{HaapasaloThesis}} Quotes are
hard to make neatly since you should use only as much as needed
without changing the text. Moreover, you often do not really
understand what the author has mentioned with his wordings if you
cannot write the same with your own words. Remember also that never
cut and paste anything without marking the quotation marks right away,
and in general, never cut and paste anything at all!

Sometimes getting the original source can be almost impossible. In an
extremely desperate situation, you can refer with structure \emph{mr
  X~[\ldots] according to ms Y~[\ldots] defined that}, if you find a
source that refers to the original source. Note also that the
reference marking is never used as sentence element (example of how 
\textbf{not} to do it: \emph{\cite{HaapasaloThesis} describes
an optimal algorithm for indexing multiversiond databases.}).



\chapter{Research Question}
\label{chapter:research question}

\fixme{\begin{itemize}
\item The threat analysis and threat assessment flows that are done and should be done in industry contexts, and with regards to OneID in particular.
\item What security flaws happen and what are especially dangerous in identity management systems, especially in the context of the human factor.
\item Add evaluation criteria!
\end{itemize}}


\chapter{Analysis}
\label{chapter:analysis}


 
\chapter{Evaluation}
\label{chapter:evaluation}

\fixme{How well did it go gets put here.}
 

  
\chapter{Conclusions}
\label{chapter:conclusions}

\fixme{Wrap-up here.
Basically the whole thing in a nutshell.
Written at the end.}


% Load the bibliographic references
% ------------------------------------------------------------------
% You can use several .bib files:
% \bibliography{thesis_sources,ietf_sources}
\bibliography{sources}


% Appendices go here
% ------------------------------------------------------------------
% If you do not have appendices, comment out the following lines
\appendix
\chapter{First appendix}
\label{chapter:first-appendix}

This is the first appendix. You could put some test images or verbose data in an
appendix, if there is too much data to fit in the actual text nicely.

For now, the Aalto logo variants are shown in Figure~\ref{fig:aaltologo}.

\begin{figure}
\begin{center}
\subfigure[In English]{\includegraphics[width=.8\textwidth]{images/aalto-logo-en}}
\subfigure[Suomeksi]{\includegraphics[width=.8\textwidth]{images/aalto-logo-fi}}
\subfigure[P� svenska]{\includegraphics[width=.8\textwidth]{images/aalto-logo-se}}
\caption{Aalto logo variants}
\label{fig:aaltologo}
\end{center}
\end{figure}



% End of document!
% ------------------------------------------------------------------
% The LastPage package automatically places a label on the last page.
% That works better than placing a label here manually, because the label might not go to the actual last page, if LaTeX needs to place floats (that is, figures, tables, and such) to the end of the document.
\end{document}
